\documentclass[12pt, a4paper]{article}

% ─── Pacotes ───────────────────────────────────────────────────────────────────
\usepackage[utf8]{inputenc}
\usepackage[T1]{fontenc}
\usepackage[brazil]{babel}
\usepackage{geometry}
\usepackage{setspace}
\usepackage{indentfirst}
\usepackage{amsmath, amssymb}
\usepackage{graphicx}
\usepackage{booktabs}
\usepackage{multirow}
\usepackage{hyperref}
\usepackage{cite}
\usepackage{xcolor}
\usepackage{caption}
\usepackage{float}
\usepackage{enumitem}
\usepackage{array}
\usepackage{makecell}

% ─── Layout ────────────────────────────────────────────────────────────────────
\geometry{
  top=3cm, bottom=2cm,
  left=3cm, right=2cm
}
\setstretch{1.5}
\setlength{\parindent}{1.25cm}

% ─── Metadados do hyperref ──────────────────────────────────────────────────────
\hypersetup{
  colorlinks=true,
  linkcolor=black,
  citecolor=black,
  urlcolor=blue,
  pdftitle={Benchmark de Modelos Preditivos para Séries Temporais Curtas},
  pdfauthor={Kailane Eduarda Felix da Silva},
}

% ═══════════════════════════════════════════════════════════════════════════════
\begin{document}

% ─── Capa ──────────────────────────────────────────────────────────────────────
\begin{titlepage}
  \centering
  \vspace*{1cm}

  {\large \textbf{UNIVERSIDADE FEDERAL DE PERNAMBUCO}} \\[0.3cm]
  {\large Centro de Informática} \\[0.3cm]
  {\large Bacharelado em Ciência da Computação} \\

  \vspace{3cm}

  {\Large \textbf{Benchmark de Modelos Preditivos para Séries Temporais Curtas: um estudo aplicado a programas rurais em Pernambuco}} \\

  \vspace{3cm}

  {\large Kailane Eduarda Felix da Silva} \\

  \vfill

  {\large Recife, \the\year}
\end{titlepage}

% ─── Resumo ────────────────────────────────────────────────────────────────────
\newpage
\begin{abstract}
\noindent

\end{abstract}

\vspace{0.5cm}
\noindent\textbf{Palavras-chave:} séries temporais, previsão de demanda,
ARIMA, Random Forest, LightGBM, walk-forward cross-validation,
políticas públicas.

% ─── Sumário ───────────────────────────────────────────────────────────────────
\newpage
\tableofcontents

% ═══════════════════════════════════════════════════════════════════════════════
\newpage
\section{Introdução}
\label{sec:introducao}

A capacidade de prever com antecedência o número de beneficiários de programas
de apoio rural é fundamental para o planejamento e alocação eficiente de recursos
públicos. No entanto, dados históricos em granularidade municipal costumam
abranger poucos anos, impondo restrições severas ao treinamento de modelos de
aprendizado de máquina.

Este trabalho propõe um benchmark sistemático de modelos preditivos aplicados
a séries temporais anuais de beneficiários de três programas rurais em
Pernambuco --- Fruticultura Irrigada, Zona Canavieira e Pesca Artesanal ---,
cobrindo o período de 2021 a 2025 e 101 municípios. O objetivo central é
identificar qual família de modelos oferece a melhor relação entre acurácia e
robustez quando o histórico disponível é limitado a cinco pontos por série.

\subsection{Motivação}

\subsection{Objetivos}

Os objetivos específicos deste trabalho são:

\begin{itemize}[leftmargin=2cm]
  \item Comparar modelos estatísticos (ARIMA, ETS), de aprendizado de máquina
        clássico (Regressão Linear, Árvore de Decisão) e moderno
        (Random Forest, LightGBM) numa tarefa de previsão com séries curtas;
  \item Avaliar o impacto de diferentes estratégias de treinamento ---
        modelo global versus modelo por setor (\textit{POR\_SETOR});
  \item Implementar validação \textit{walk-forward} com múltiplos folds para
        medir a estabilidade dos modelos, evitando conclusões baseadas em
        um único ano de teste;
  \item Identificar e documentar fontes de \textit{data leakage} comuns em
        séries temporais curtas e suas consequências nos resultados reportados.
\end{itemize}

\subsection{Organização do Texto}

O restante do trabalho está organizado da seguinte forma:
a Seção~\ref{sec:referencial} apresenta o referencial teórico;
a Seção~\ref{sec:dados} descreve os dados e o pré-processamento;
a Seção~\ref{sec:metodologia} detalha os modelos e o protocolo de validação;
a Seção~\ref{sec:resultados} apresenta e discute os resultados;
a Seção~\ref{sec:conclusao} conclui o trabalho e aponta direções futuras.

% ═══════════════════════════════════════════════════════════════════════════════
\newpage
\section{Referencial Teórico}
\label{sec:referencial}

% Sugestão de subseções — preencha com ~1 página cada:

\subsection{Previsão de Séries Temporais Curtas}

\subsection{Modelos Estatísticos Clássicos}

\subsubsection{ARIMA}

\subsubsection{Suavização Exponencial (ETS)}

\subsection{Aprendizado de Máquina para Séries Temporais}

\subsubsection{Engenharia de Features com Lags}


\subsubsection{Random Forest}


\subsubsection{LightGBM}


\subsection{Validação de Modelos Temporais}


\subsection{Data Leakage em Séries Temporais}


% ═══════════════════════════════════════════════════════════════════════════════
\newpage
\section{Dados e Pré-processamento}
\label{sec:dados}

\subsection{Descrição do Conjunto de Dados}

Os dados utilizados neste trabalho são registros anuais de beneficiários de
três programas rurais administrados pelo [nome da instituição/secretaria],
cobrindo 101 municípios de Pernambuco entre 2021 e 2025. O conjunto original
contém 517 observações distribuídas entre três programas: Fruticultura Irrigada,
Zona Canavieira e Pesca Artesanal.

\begin{table}[H]
  \centering
  \caption{Resumo do conjunto de dados}
  \label{tab:dados_resumo}
  \begin{tabular}{ll}
    \toprule
    \textbf{Atributo}          & \textbf{Valor}              \\
    \midrule
    Período                    & 2021–2025 (5 anos)          \\
    Municípios                 & 101                         \\
    Programas                  & 3                           \\
    Observações originais      & 517                         \\
    Granularidade              & Município $\times$ Programa $\times$ Ano \\
    \bottomrule
  \end{tabular}
\end{table}

Uma característica relevante do conjunto é que nem todo município possui
registros para todos os programas em todos os anos, gerando lacunas na série.
Essas lacunas são tratadas conforme descrito na Seção~\ref{sec:grid}.

\subsection{Pré-processamento}

\subsubsection{Limpeza e Padronização}


\subsubsection{Expansão do Grid Completo}
\label{sec:grid}

Para garantir que todos os modelos operem sobre séries sem lacunas temporais,
foi criado o produto cartesiano completo de municípios, programas e anos
já presentes no conjunto original. Combinações sem registro real receberam
o valor zero de beneficiários, estratégia conservadora que assume ausência
de atividade em vez de dado faltante.

Após a expansão, o conjunto passou de 517 para X observações.

\subsubsection{Codificação de Variáveis Categóricas}


\subsection{Engenharia de Features}
\label{sec:features}

Todas as features foram construídas de forma a não introduzir
\textit{data leakage} --- ou seja, nenhuma delas utiliza informações
do período de teste nem do ano imediatamente anterior ($t-1$) ao alvo,
conforme discutido na Seção~\ref{sec:leakage}.

\begin{table}[H]
  \centering
  \caption{Features utilizadas nos modelos de ML (versão final, v5)}
  \label{tab:features}
  \begin{tabular}{p{3.5cm} p{8.5cm}}
    \toprule
    \textbf{Feature}         & \textbf{Descrição} \\
    \midrule
    \texttt{lag\_2}          & Beneficiários em $t-2$ \\
    \texttt{lag\_3}          & Beneficiários em $t-3$ \\
    \texttt{rolling\_mean\_2}& Média de \texttt{lag\_2} e \texttt{lag\_3}: $(lag_2 + lag_3)/2$ \\
    \texttt{trend}           & Variação absoluta: $lag_2 - lag_3$ \\
    \texttt{growth\_rate}    & Variação proporcional: $(lag_2 - lag_3)/lag_3$; zero quando $lag_3 = 0$ \\
    \texttt{zero\_historico} & Indicador binário: 1 se $lag_2 = lag_3 = 0$ \\
    \texttt{ano\_rel}        & Anos desde o início da série (tendência linear) \\
    \texttt{programa\_total\_lag2} & Soma de \texttt{lag\_2} de todos os municípios do programa no mesmo ano \\
    \texttt{share\_municipio}& Peso relativo do município no programa: $lag_2 / (programa\_total\_lag_2 + 1)$ \\
    \bottomrule
  \end{tabular}
\end{table}

% ═══════════════════════════════════════════════════════════════════════════════
\newpage
\section{Metodologia}
\label{sec:metodologia}

\subsection{Modelos Avaliados}

Seis modelos foram selecionados, representando três famílias distintas,
com o objetivo de cobrir um espectro amplo de complexidade e suposições
sobre os dados.

\begin{table}[H]
  \centering
  \caption{Modelos avaliados por família}
  \label{tab:modelos}
  \begin{tabular}{lll}
    \toprule
    \textbf{Família}   & \textbf{Modelo}         & \textbf{Configuração} \\
    \midrule
    Estatístico        & ARIMA                   & Ordem $(1,1,0)$ \\
    Estatístico        & ETS (Holt Linear)       & Tendência aditiva, sem sazonalidade \\
    ML Clássico        & Regressão Linear        & Sem regularização \\
    ML Clássico        & Árvore de Decisão       & \texttt{max\_depth=4} \\
    ML baseado em Ensemble         & Random Forest           & 200 estimadores, \texttt{max\_depth=6}, \texttt{min\_samples\_leaf=3} \\
    ML baseado em Ensemble         & LightGBM                & 200 estimadores, \texttt{lr=0.05}, \texttt{max\_depth=4} \\
    \bottomrule
  \end{tabular}
\end{table}

\subsubsection{Modelos Estatísticos}


\subsubsection{Modelos de ML}

\subsection{Modos de Treinamento}

Cada modelo foi avaliado em dois modos:

\begin{itemize}[leftmargin=2cm]
  \item \textbf{GLOBAL:} um único modelo treinado com dados de todos os
        municípios e programas simultaneamente.
  \item \textbf{POR\_SETOR:} um modelo treinado separadamente para cada
        programa, permitindo que o modelo capture dinâmicas específicas
        de cada setor.
\end{itemize}

\subsection{Métrica de Avaliação}

A métrica principal adotada é o Erro Percentual Absoluto Médio
(\textit{Mean Absolute Percentage Error}, MAPE), definido como:

\begin{equation}
  \text{MAPE} = \frac{1}{n} \sum_{t=1}^{n}
  \left| \frac{y_t - \hat{y}_t}{y_t} \right| \times 100\%
  \label{eq:mape}
\end{equation}

\noindent onde $y_t$ é o valor real e $\hat{y}_t$ é o valor previsto no
período $t$. Observações com $y_t \approx 0$ são excluídas do cálculo
para evitar distorções numéricas.

Para garantir comparabilidade entre modelos estatísticos (que operam em
séries agregadas por programa) e modelos de ML (que operam no nível
município), as previsões de ML são somadas por programa antes do cálculo
do MAPE. Essa função é denominada \texttt{evaluate\_aggregated()} na
implementação.

São reportados também MAE e RMSE como métricas complementares.

\subsection{Protocolo de Validação}
\label{sec:validacao}

\subsubsection{Benchmark de Fold Único}


\subsubsection{Walk-forward Cross-validation}

A validação principal utiliza a estratégia \textit{expanding window}
(janela expansível), em que o conjunto de treino cresce a cada fold
enquanto o conjunto de teste avança um ano por vez. Com dados de 2021
a 2025 e mínimo de dois anos de treino, foram gerados três folds:

\begin{table}[H]
  \centering
  \caption{Folds do walk-forward cross-validation}
  \label{tab:folds}
  \begin{tabular}{ccc}
    \toprule
    \textbf{Fold} & \textbf{Treino}   & \textbf{Teste} \\
    \midrule
    1             & 2021–2022         & 2023           \\
    2             & 2021–2023         & 2024           \\
    3             & 2021–2024         & 2025           \\
    \bottomrule
  \end{tabular}
\end{table}

O desempenho final de cada modelo no CV é reportado como média e desvio
padrão do MAPE entre os folds válidos (modelos estatísticos requerem
no mínimo 3 pontos de treino; folds com série insuficiente são excluídos
automaticamente).

\subsection{Problemas Identificados e Soluções}
\label{sec:leakage}

\subsubsection{Problema 1 — Comparação em Granularidades Distintas}

\subsubsection{Problema 2 — Data Leakage via \texttt{lag\_1}}

A inclusão de \texttt{lag\_1} (valor de $t-1$) como feature de ML
equivale, em séries estáveis de curto horizonte, a fornecer ao modelo
a resposta que ele deveria prever. Isso ficou evidenciado pelo desempenho
da Árvore de Decisão na Pesca Artesanal: MAPE de 0,03\% com
\texttt{lag\_1} respondendo por 98,4\% da importância de features ---
o modelo simplesmente repetia o ano anterior.

\begin{table}[H]
  \centering
  \caption{Evidência do leakage via \texttt{lag\_1} — Árvore de Decisão, Pesca Artesanal}
  \label{tab:leakage_lag1}
  \begin{tabular}{lrrrr}
    \toprule
    \textbf{Município} & \textbf{Real 2025} & \textbf{Previsto} & \textbf{Erro abs.} & \textbf{lag\_1} \\
    \midrule
    Goiana    & 1817 & 1810,00 &  7,00 & 1810 \\
    Recife    &  600 &  596,00 &  4,00 &  592 \\
    Ipojuca   &  186 &   89,31 & 96,69 &   98 \\
    Igarassu  &  511 &  443,25 & 67,75 &  477 \\
    \bottomrule
  \end{tabular}
\end{table}

\noindent\textbf{Solução:} \texttt{lag\_1} foi removido inteiramente do
pré-processamento. As features \texttt{rolling\_mean\_2} e \texttt{trend},
que dependiam internamente de \texttt{lag\_1}, foram recalculadas usando
\texttt{lag\_2} e \texttt{lag\_3}.

\subsubsection{Problema 3 — Leakage Residual em Features Derivadas}

% ═══════════════════════════════════════════════════════════════════════════════
\newpage
\section{Resultados e Discussão}
\label{sec:resultados}

\subsection{Benchmark de Fold Único (2025)}


\begin{table}[H]
  \centering
  \caption{Resultados do benchmark — fold único (ano de teste: 2025)}
  \label{tab:resultados_benchmark}
  \begin{tabular}{llll}
    \toprule
    \textbf{Modelo}       & \textbf{Modo}    & \textbf{Programa}          & \textbf{MAPE\%} \\
    \midrule
    RandomForest          & POR\_SETOR       & Fruticultura Irrigada      & 2,30  \\
    RandomForest          & POR\_SETOR       & Zona Canavieira            & 3,00  \\
    LightGBM              & POR\_SETOR       & Fruticultura Irrigada      & 3,30  \\
    RandomForest          & GLOBAL           & Todos                      & 3,76  \\
    ETS (Holt Linear)     & POR\_SETOR       & Zona Canavieira            & 3,78  \\
    ARIMA                 & GLOBAL           & Todos                      & 4,06  \\
    ARIMA                 & POR\_SETOR       & Fruticultura Irrigada      & 4,31  \\
    DecisionTree          & POR\_SETOR       & Pesca Artesanal            & 7,48  \\
    RandomForest          & POR\_SETOR       & Pesca Artesanal            & 7,74  \\
    ARIMA                 & POR\_SETOR       & Pesca Artesanal            & 9,32  \\
    \bottomrule
  \end{tabular}
\end{table}

\subsection{Walk-forward Cross-validation}

\begin{table}[H]
  \centering
  \caption{Walk-forward CV — média de 3 folds (2023, 2024, 2025)}
  \label{tab:resultados_cv}
  \begin{tabular}{lllrrl}
    \toprule
    \textbf{Modelo}   & \textbf{Modo} & \textbf{Programa}       & \textbf{MAPE médio} & \textbf{Desvio} & \textbf{Folds} \\
    \midrule
    ARIMA             & POR\_SETOR    & Fruticultura Irrigada   & 5,84\%              & ±2,16           & 2 \\
    ARIMA             & POR\_SETOR    & Pesca Artesanal         & 6,75\%              & ±3,64           & 2 \\
    ARIMA             & GLOBAL        & Todos                   & 7,46\%              & ±4,81           & 2 \\
    RandomForest      & POR\_SETOR    & Fruticultura Irrigada   & 7,83\%              & ±9,78           & 3 \\
    DecisionTree      & POR\_SETOR    & Pesca Artesanal         & 9,55\%              & ±7,68           & 3 \\
    RandomForest      & POR\_SETOR    & Pesca Artesanal         & 9,91\%              & ±5,70           & 3 \\
    DecisionTree      & POR\_SETOR    & Fruticultura Irrigada   & 10,51\%             & ±2,80           & 3 \\
    \bottomrule
  \end{tabular}
\end{table}

\subsection{Discussão}

\subsubsection{POR\_SETOR versus GLOBAL}

\subsubsection{ARIMA como Modelo Mais Robusto}

\subsubsection{Alta Variância dos Modelos de ML no CV}


\subsubsection{Limitações}

\begin{itemize}[leftmargin=2cm]
  \item \textbf{Tamanho da série:} cinco pontos anuais por série limitam
        a generalização de qualquer modelo e impõem incerteza elevada nas
        estimativas de CV.
  \item \textbf{Zeros estruturais:} municípios sem histórico em um dado
        programa recebem valor zero, o que pode distorcer features baseadas
        em lags e inflacionar o MAPE quando o valor real também é zero.
  \item \textbf{Ausência de sazonalidade:} com dados anuais, não é
        possível modelar padrões intra-anuais que possam existir nos programas.
\end{itemize}

% ═══════════════════════════════════════════════════════════════════════════════
\newpage
\section{Conclusão}
\label{sec:conclusao}

% ═══════════════════════════════════════════════════════════════════════════════
\newpage
\bibliographystyle{IEEEtran}
\begin{thebibliography}{00}

\bibitem{hyndman2021}
R. J. Hyndman e G. Athanasopoulos,
\textit{Forecasting: Principles and Practice}, 3ª ed.
OTexts, 2021. Disponível em: \url{https://otexts.com/fpp3/}

\bibitem{makridakis2022}
S. Makridakis, E. Spiliotis e V. Assimakopoulos,
``M5 accuracy competition: Results, findings, and conclusions,''
\textit{International Journal of Forecasting}, vol. 38, n. 4,
pp. 1346--1364, 2022.

\bibitem{breiman2001}
L. Breiman,
``Random forests,''
\textit{Machine Learning}, vol. 45, n. 1, pp. 5--32, 2001.

\bibitem{ke2017}
G. Ke et al.,
``LightGBM: A highly efficient gradient boosting decision tree,''
\textit{Advances in Neural Information Processing Systems (NeurIPS)},
vol. 30, 2017.

\bibitem{box2015}
G. E. P. Box, G. M. Jenkins, G. C. Reinsel e G. M. Ljung,
\textit{Time Series Analysis: Forecasting and Control}, 5ª ed.
Wiley, 2015.

\bibitem{bergmeir2012}
C. Bergmeir e J. M. Benítez,
``On the use of cross-validation for time series predictor evaluation,''
\textit{Information Sciences}, vol. 191, pp. 192--213, 2012.

\bibitem{sklearn}
F. Pedregosa et al.,
``Scikit-learn: Machine learning in Python,''
\textit{Journal of Machine Learning Research}, vol. 12,
pp. 2825--2830, 2011.

\end{thebibliography}

\end{document}